% !TeX TS-program = xelatex
% !TeX program = xelatex
% !TeX encoding = UTF-8
% NC-style two-column working manuscript. Nature Communications does not
% require a dedicated two-column class at initial submission; this layout is
% for internal review and visual preparation only.
\documentclass[10pt,a4paper,twocolumn]{article}

% Nature Communications initial-submission manuscript framework.
% Compile this Chinese working file with XeLaTeX.
\usepackage[
  top=1.8cm,
  bottom=1.8cm,
  left=1.65cm,
  right=1.65cm,
  columnsep=0.72cm
]{geometry}
\usepackage{fontspec}
\usepackage{xeCJK}
\setmainfont{Times New Roman}
\setCJKmainfont[AutoFakeBold=2]{SimSun}
\setCJKsansfont{SimHei}
\setCJKmonofont{SimSun}

\usepackage{amsmath,amssymb,bm}
\usepackage{graphicx}
\usepackage{setspace}
\usepackage[numbers]{natbib}
\usepackage[hidelinks]{hyperref}

\graphicspath{{figures/}}
\setstretch{1.08}
\setlength{\parindent}{2em}
\setlength{\parskip}{0pt}
\raggedbottom

\title{交通风险传播中的信息介导非最近邻效应}
\author{禹尧珅}
\date{}

\begin{document}

\maketitle

\begin{abstract}
% 摘要在主要结果冻结后写入。
\end{abstract}

\section*{Introduction}
道路交通是由连续物理交互与日益密集的信息连接共同构成的复杂系统。局部制动、速度下降和瓶颈激活可以经由跟驰关系、车辆运动和交通波在车辆串与路网中传播，使原本局部的状态扰动产生跨越时间和空间的影响\cite{saberi2020contagion}\cite{duan2023bottlenecks}。联网与自动驾驶技术进一步改变了这种传播结构：车辆不仅响应邻近交通状态，还能够通过车际通信、路侧感知或云端控制获取并采用来自非邻近位置的信息\cite{wang2025delays}。交通风险因而不再只受单一物理交互网络塑造，而可能由物理传播与信息传播在不同拓扑和时间尺度上的共同作用所决定；厘清这两类传播如何分别改变车辆动作和风险结果，是理解联网交通系统安全机制的基础。

现有研究从风险时空演化、非局部交通流建模、逐跳风险级联和联网车辆控制等方向揭示了交通状态跨空间作用的不同侧面。基于轨迹和替代安全指标的研究将冲突易发位置组织为图结构，用于刻画交叉口风险的时空依赖\cite{wang2024conflict}；风险场驾驶模型则把感知风险转化为速度和横向位置选择\cite{kolekar2020risk}。非局部交通流模型将前方距离、空间邻域密度或多车状态纳入运动规律，用以分析前视信息对交通波和均匀流稳定性的影响\cite{sun2020lookahead}\cite{huang2022nonlocal}。

网络扰动还受到传播路径、节点度和基本模体制约，特定结构能够隔离故障扩散；交通基础设施中的需求分布和瓶颈结构也会改变网络响应\cite{bao2022motifs}\cite{kaiser2021isolators}\cite{hamedmoghadam2021percolation}。车队级联研究沿连续物理邻接关系递推临界减速度或冲突状态，解释风险如何经中间车辆累积、放大或衰减\cite{luo2026crm}。联网车辆控制研究进一步考察通信资源、时延、失效和动态拓扑对闭环稳定性与安全性的影响\cite{razzaghpour2023controlaware}\cite{wang2024topology}；碰撞规避建模、密集学习和对抗安全测试则扩展了复杂行为与稀有风险的分析范围\cite{schumann2026collision,feng2026dense,xu2026ga2ad}。这些工作尚未直接回答一条来源可追溯的远端消息是否被目标车辆采用，并由此改变目标区域的总体风险。

这些研究尚不能单独识别指定远端信息造成的区域级风险效应。风险图描述的是风险状态之间的时空依赖，而不是对信息通道的干预；非局部交通流模型把空间邻域状态直接写入运动规律，却不区分离散消息是否生成、发送、交付并被控制器采用；通信控制模型虽然把接收信息接入车辆控制，但信息图中存在连接、消息到达目标车辆和控制器实际使用该消息仍是不同事件。更重要的是，空间上不相邻的源车辆与目标车辆之间仍可能存在多跳跟驰、交通波或车辆运动构成的物理风险传播链，共同交通状态和共享控制规则也可能使两端同步变化。因此，本文要解决的不是远端通信是否一般性地改善控制性能，而是能否在排除物理传播链和共同状态等替代解释后，识别一条来源可追溯的风险消息在源扰动的物理影响到达前，通过改变目标车辆动作而对目标区域总体风险产生可归因效应。

本文将来源可追溯的远端风险消息被目标区域内符合资格的目标车辆采用后，在源扰动的物理影响到达前引起目标区域总体风险变化的现象，定义为信息介导的非最近邻风险效应，简称“类超距作用”。\footnote{“类超距作用”仅用于概括远隔区域之间由显式信息介导的风险效应，不表示量子纠缠、超光速传播或无媒介作用。}

源区域与目标区域分别记为道路状态空间中的 $\mathcal R^{\mathsf S}$ 和 $\mathcal R^{\mathsf T}$。前者定义源风险事件的发生位置，后者定义目标车辆集合及其区域风险结果的观测范围，且 $\mathcal R^{\mathsf S}\cap\mathcal R^{\mathsf T}=\varnothing$。对每个配对单元 $\omega$，在处理开始前依据冻结的区域、时间和车辆资格规则确定目标车辆集合 $\mathcal J^{\mathsf T}_{\omega}$；单车动作和风险量是该集合的组成观测，目标区域总体风险由集合内单车风险聚合得到。区域归属依据道路边、车道、纵向位置和路线组件判定，不以欧氏距离单独判定。区域不重叠也不表示逐跳物理路径必然不存在，后者必须通过时序可达性审计确认。本文要求证据链依次覆盖源风险、消息生成与发送、交付与校验、控制采用、实际动作和区域风险变化，并以屏蔽指定源消息的配对条件估计其增量效应。

本文据此构建一个区分物理交通网络与显式信息网络的可证伪识别框架，并以消息生命周期和四格反事实干预估计指定远端信息对目标区域总体风险的增量效应。证据按照识别强度与场景现实性逐层推进：首先在两条物理断开的走廊中隔离交通波、多跳跟驰和车辆运动，检验指定信息是否产生区域级增量风险效应；随后审计消息从生成到风险变化的完整生命周期，排除错误来源、过期信息和仅交付未采用等替代解释；继而将冻结机制迁移到物理连通的合流路网，比较信息响应与物理影响的到达顺序，并评估闭环控制价值；再利用真实轨迹与通信数据校核风险测量和消息可观测层级；最后通过改变控制器、风险类型、信息图和通信条件界定适用范围。这个递进设计将逐跳物理传播作为必要的竞争基线，并把本文的中心问题限定为：在什么可观测条件和因果边界内，远端风险信息能够通过目标车辆动作改变目标区域的总体风险。


\section*{Results}

\subsection*{Result 1：测量与分析链的有效性}

本节判断测量、运行和分析链能否可靠地产生后续结果所需的证据。它是全部实证分析的质量门槛，只回答研究工具是否有效，不回答核心效应是否存在。本节报告校准表现、复现性、判定准确性及未通过项目，不解释主要处理效应，也不以校准通过替代因果证据。

\subsection*{Result 2：非最近邻风险效应的核心识别}

本节回答在排除源端至目标端物理传播的条件下，开放指定消息通道是否对目标区域总体风险产生可归因的增量效应。主要效应使用全部预先登记的正式四格单元，消息未采用或单车动作未分化的运行也保留在意向运行分母。本节报告区域效应方向、实际意义、不确定性、工程失败和缺失数据界限，不解释消息作用链，不评价物理连通场景中的控制价值，也不主张跨条件推广。

\subsection*{Result 3：信息—动作—风险作用链的机制审计}

本节判断 Result 2 的风险差异能否沿来源、传输、校验、采用、实际输入和风险的顺序追溯，并通过负对照排除关键替代解释。它报告作用链通过率、时间顺序和失败环节，但机制是否通过不改变 Result 2 的样本成员关系。本节不重新估计核心主效应，也不把消息到达、采用次数或日志完整性单独解释为因果证据。

\subsection*{Result 4：物理连通场景中的信息先行与控制价值}

本节回答在物理传播路径恢复后，信息响应能否先于物理影响到达，并在安全约束下形成有实际价值的闭环控制结果。它负责从结构隔离的识别场景迁移到更现实的交通交互场景。本节报告信息与物理到达的时间关系、风险与控制收益、安全约束、效率代价和失败运行；这些结果不能替代 Result 2 对物理路径的结构隔离，也不能以风险下降单独证明控制有效。

\subsection*{Result 5：真实数据支持的外部校核与证据边界}

本节判断真实轨迹与通信数据能够校核哪些测量、分布和可观测性，并明确各数据源实际支持的证据层级。它负责外部一致性检查，而不是把不同数据源拼接为现实闭环因果证据。本节报告可支持的校核结果、不可观测字段、数据质量问题和结论限制，不从缺少消息采用或动作响应字段的数据推断真实控制效应，也不直接比较定义不同的风险量。

\subsection*{Result 6：适用域、稳健性与失效边界}

本节回答前述证据在控制器、风险类型、交通拓扑、通信条件和尾部场景变化时能否保持，并据此界定机制的适用域与失效域。它承担一般性和边界条件的综合判断。本节报告跨条件效应、异质性、稳健性、安全与效率权衡，以及零效应、不确定和失效条件；不重复 Methods 中的条件定义，不重复单一基准场景的主要效应，也不将有限扩展表述为普遍规律。

\section*{Discussion}
本节解释主要发现的机制含义、适用边界、与既有研究的关系以及对联网交通安全设计的影响，并区分由仿真、真实数据校核和不可观测字段分别支持的结论。

\section*{Methods}
\subsection*{System model and two-stage framework}
本文研究由车辆运动、物理交互和显式消息传递共同构成的联网交通系统。该系统既产生沿道路和车辆邻接关系传播的物理影响，也允许远端风险信息经通信和控制器采用改变车辆动作。设 $\mathcal V=\{1,\ldots,N\}$ 为车辆集合，$X_t=\{x_{i,t}:i\in\mathcal V\}$ 为全体车辆状态，$U_t=\{u_{i,t}:i\in\mathcal V\}$ 为实际控制输入，$G^{\mathsf P}(t)$ 和 $G^{\mathsf I}(t)$ 分别表示物理交通图和显式信息图，$M_t$ 表示截至时刻 $t$ 已生成且尚在协议生命周期内的消息集合。研究的共同系统模型写为
\begin{equation}
\Sigma_t=\left(X_t,U_t,G^{\mathsf P}(t),G^{\mathsf I}(t),M_t\right),
\end{equation}
该模型是两个研究阶段的共同基础：Stage I 在给定系统和控制逻辑下识别指定消息的增量因果效应；仅当该作用达到预设证据门槛后，Stage II 才优化消息采用和车辆动作。两阶段共享状态、图、消息和风险测量，但具有不同的问题、约束和评价标准。

车辆类型记为 $\kappa_i\in\{\mathsf{AV},\mathsf{HV}\}$。$\mathsf{AV}$ 和 $\mathsf{HV}$ 只表示控制器类别，不改变车辆索引。对每个时刻 $t$，$\mathcal{N}_i^{\mathsf P}(t)$ 表示能够通过跟驰、合流或冲突直接影响车辆 $i$ 的物理邻居，$\mathcal{N}_i^{\mathsf I}(t)$ 表示通信协议允许向车辆 $i$ 提供消息的来源集合。由有向序对 $(k,i)$ 表示从车辆 $k$ 指向车辆 $i$，物理图和信息图定义为
\begin{equation}
\begin{aligned}
\mathcal{E}^{\mathsf P}(t)&=\{(k,i):k\in\mathcal{N}_i^{\mathsf P}(t),\ k\neq i\},\\
\mathcal{E}^{\mathsf I}(t)&=\{(k,i):k\in\mathcal{N}_i^{\mathsf I}(t),\ k\neq i\},\\
G^{\mathsf P}(t)&=\bigl(\mathcal{V},\mathcal{E}^{\mathsf P}(t)\bigr),\qquad
G^{\mathsf I}(t)=\bigl(\mathcal{V},\mathcal{E}^{\mathsf I}(t)\bigr).
\end{aligned}
\end{equation}
物理边表示车辆运动上的直接作用；沿连续物理边形成的逐跳链属于物理传播。信息边只表示协议允许通信，不表示消息已经生成、发送、交付或被控制器采用。两类图可以同时存在，也可以在双走廊设计中被结构性分离；是否存在物理传播不能只由空间距离或某一时刻没有直接边来判断。

车辆 $i$ 在时刻 $t$ 的状态写为
\begin{equation}
x_{i,t}=\bigl(\mathbf p_{i,t},v_{i,t},a_{i,t},\ell_{i,t}\bigr),
\end{equation}
其中 $\mathbf p_{i,t}$ 为二维位置，$v_{i,t}$ 为速度，$a_{i,t}$ 为实际纵向加速度，$\ell_{i,t}$ 为车道和道路边状态。$u_{i,t}^{\mathsf{cmd}}$ 表示控制器请求的纵向或横向命令，$u_{i,t}$ 表示该命令经过车辆执行器和安全约束后实际施加的控制输入，$\xi_{i,t}$ 表示外生随机输入。离散仿真时钟上的状态更新写为
\begin{equation}
x_{i,t+1}=f_i^{\mathsf{dyn}}\!\left(x_{i,t},u_{i,t},\xi_{i,t};G^{\mathsf P}(t)\right).
\end{equation}
函数 $f_i^{\mathsf{dyn}}$ 包含车辆运动学、道路边界和物理邻居作用；其具体跟驰、制动、换道和合流规则由实验配置指定，不在本节预设某一种控制器。所有车辆状态、控制输入和消息事件使用同一仿真时钟记录。

前述消息集合 $M_t$ 中的每条消息由唯一标识 $q$、来源车辆 $i$、目标车辆 $j$、消息载荷和协议时间戳组成，记为 $m_{q,i,j,t}$。协议允许发往目标车辆 $j$ 的候选消息集合为
\begin{equation}
M_{j,t}^{\mathsf{cand}}=\left\{m_{q,i,j,t}\in M_t:
(i,j)\in\mathcal{E}^{\mathsf I}(t)\right\}.
\end{equation}
候选集合表示消息已经生成且协议允许该来源向目标发送，但不表示消息已经发送、交付或可供控制器使用。消息的生成、发送、交付、字段校验、采用和失效在日志中分别记录；这些事件的先后和采用规则在 Stage I 中定义。信息边不绕过目标车辆的本地感知和安全约束，也不自动改变车辆动作。

系统模型沿用前文定义的源区域 $\mathcal R^{\mathsf S}$、目标区域 $\mathcal R^{\mathsf T}$ 和目标车辆集合 $\mathcal J^{\mathsf T}_{\omega}$。区域归属用于确定源风险事件和目标区域风险结果的观测范围；道路边、车道、纵向坐标范围、路线组件、拓扑距离下限及车辆资格规则，在后文的 Experimental design and data sources 中按场景给出并在确认性实验前冻结。源风险、远端消息、目标车辆动作和区域风险结果由此形成“源事件—信息传递—控制采用—区域响应”的共同分析链，物理图则给出与该链竞争的普通交通传播路径。

源车辆 $i$ 和目标车辆 $j\in\mathcal J^{\mathsf T}_{\omega}$ 的资格由区域和时刻共同定义：源车辆的合格风险事件必须在 $\mathcal R^{\mathsf S}$ 内触发，目标车辆集合由处理前规则确定，并在结果窗口内作为目标区域风险的组成单位。车辆在运行过程中离开区域不自动改变已记录的身份，但若离开导致目标状态、风险窗口或路径审计无法满足预设规则，则该配对按失败或不确定处理。物理传播路径按时序可达性审计：若扰动可沿物理边逐跳到达 $j$，记录最早物理到达时刻 $\tau_j^{\mathsf P}$；若不存在可达路径，则令 $\tau_j^{\mathsf P}=+\infty$；若观察期内无法判定，则保留为不确定或右删失，不当作路径不存在。

仿真采用统一的离散时钟 $t=n\Delta t$，其中 $n$ 为时间步索引，$\Delta t$ 在数值收敛校准后固定。每个时间步依次读取车辆状态和物理邻接关系，判定源风险事件，更新消息的生成、发送、交付、校验和采用状态，计算控制输入，并由交通仿真器推进车辆运动。风险量由同一时钟下的状态和交互关系计算。状态、消息、控制和风险记录共享运行标识、时间步和时间原点；事件先后依据仿真时间戳而非文件行序判断。

车辆纵向运动由交通仿真器与车辆控制器共同更新。人类驾驶车辆采用预先固定的跟驰与随机扰动模型；自动驾驶车辆采用相应实验条件下冻结的基线或控制策略。纵向命令受车辆动力学、加减速度和碰撞规避约束，横向运动由冻结的路线、换道和合流规则决定。指定远端消息只构成自动驾驶车辆控制器的附加输入，不关闭本地感知、基础控制器或安全约束。车辆类型、控制器参数和道路规则在场景运行前固定，同一四格配对内保持一致。

位置以米、速度以米每秒、加速度以米每二次方秒记录；道路边、车道、前车标识和消息标识作为离散状态保存。净间距和相对速度由同一时刻的车辆状态派生。不存在有效前车或某一字段不可观测时，该量记为缺失，不以零值替代。每辆车在单次运行中使用唯一身份，其进入、活动和离开事件按时间排序；车辆离开后不得继续被物理边、消息或风险记录引用。

物理传播采用时间展开的交通图审计。图中同时保留同一车辆相邻时间步之间的运动延续关系，以及由跟驰、合流或冲突产生的车辆间直接作用关系。一条从源车辆 $i$ 到目标车辆 $j$ 的物理路径必须沿仿真时间单调前进，且每一步均满足冻结的道路拓扑、车辆在场条件和物理边建立规则。$\tau_j^{\mathsf P}$ 取所有合格时间路径的最早到达时刻，表示物理影响可能到达的最早保守边界，而不是事后观察到的实际传播时刻。因此，要求信息响应早于 $\tau_j^{\mathsf P}$ 是保守判定。双走廊设计仅在道路拓扑断开、无跨走廊车辆转移、无共享控制状态且无合格时间路径时令 $\tau_j^{\mathsf P}=+\infty$；物理连通场景使用经独立校准的 $\tau_j^{\mathsf P}$ 保守下界。

消息载荷包含唯一消息标识、源事件标识、源车辆和目标车辆标识、风险类型、风险摘要及生成时间。每次运行分别记录车辆状态、源事件、消息生命周期、请求命令、实际输入、实际动作和风险结果。四格运行由同一冻结快照分别启动，消息缓存和运行状态互不共享。主随机种子按背景交通、人类驾驶扰动、通信和控制器等子系统派生为相互独立的命名随机流；四格使用相同的对应随机流，某一处理触发的额外随机调用不得推进其他子系统的随机序列。这样既保持共同随机数配对，也避免跨运行状态泄漏。道路边界、区域坐标、车辆进入与退出规则及场景参数在 Experimental design and data sources 中按场景给出。

\subsection*{Stage I: causal identification}
Stage I 检验本文定义的“类超距作用”：开放指定消息通道是否使来源可追溯的远端风险信息在普通物理影响到达前，通过改变目标车辆集合中的车辆动作而改变目标区域总体风险。主效应按预先登记的全部四格配对单元估计，不以消息是否采用或单车动作是否分化筛选样本；消息采用与单车动作分化用于检验作用链是否成立。该术语表示信息介导的非最近邻风险效应，相关限定见脚注。

设 $S\in\{0,1\}$ 表示指定源风险是否存在，$C\in\{0,1\}$ 表示指定消息通道是否开放。对完整四格配对单元 $\omega$，令 $R_{j,\omega}(s,c)$ 表示目标车辆 $j\in\mathcal J^{\mathsf T}_{\omega}$ 在处理条件 $(S=s,C=c)$ 下的窗口风险负担，并定义目标区域总体风险为
\begin{equation}
R_{\mathcal R^{\mathsf T},\omega}(s,c)=\sum_{j\in\mathcal J^{\mathsf T}_{\omega}}R_{j,\omega}(s,c).
\end{equation}
单元效应定义为
\begin{equation}
\begin{aligned}
\delta_{\omega}^{\mathsf{NL}}={}&[R_{\mathcal R^{\mathsf T},\omega}(1,1)-R_{\mathcal R^{\mathsf T},\omega}(1,0)]\\
&-[R_{\mathcal R^{\mathsf T},\omega}(0,1)-R_{\mathcal R^{\mathsf T},\omega}(0,0)].
\end{aligned}
\end{equation}
总体效应为 $\Delta^{\mathsf{NL}}=\mathbb E_\omega[\delta_\omega^{\mathsf{NL}}]$。这是开放指定消息通道的意向运行效应：无论某次运行中消息是否最终被采用，该运行都保留在预先登记的分析分母中。对于数值越大表示风险越高的终点，负值表示平均风险下降，正值表示风险增加。令 $\delta^{\mathsf{eq}}$ 表示可视为零效应的最大容许偏差，令 $\delta^{\mathsf R}$ 表示具有实际意义的最小风险变化幅度，并要求 $0\leq\delta^{\mathsf{eq}}<\delta^{\mathsf R}$。二者通过校准确定，并在确认性实验开始前冻结。
作为区域级辅助结果，令 $B_{\mathcal R^{\mathsf T},\omega}(s,c)\in\{0,1\}$ 表示处理条件 $(s,c)$ 下，目标区域在结果窗口内是否至少发生一次预先定义的合格风险事件，并令 $P_{\mathcal R^{\mathsf T}}(s,c)=\mathbb E_\omega[B_{\mathcal R^{\mathsf T},\omega}(s,c)]$。其对应的辅助交互效应为
\begin{equation}
\begin{aligned}
\delta_{\omega}^{\mathsf{event}}={}&[B_{\mathcal R^{\mathsf T},\omega}(1,1)-B_{\mathcal R^{\mathsf T},\omega}(1,0)]\\
&-[B_{\mathcal R^{\mathsf T},\omega}(0,1)-B_{\mathcal R^{\mathsf T},\omega}(0,0)],\\
\Delta^{\mathsf{event}}={}&\mathbb E_\omega[\delta_{\omega}^{\mathsf{event}}].
\end{aligned}
\end{equation}
该概率指标用于描述区域是否发生风险，不替代以风险负担总量为主要结果的 $R_{\mathcal R^{\mathsf T},\omega}(s,c)$。

消息处理状态由 $D_{q,j}(t)$、$Q_{q,j}(t)$ 和 $Z_{q,j}(t)$ 表示，分别对应消息 $q$ 对目标车辆 $j$ 的交付、联合校验和采用状态，三者均属于 $\{0,1\}$。令 $t_0$ 为配对试验的共同触发时刻，$t_q^{\mathsf{gen}}$、$t_q^{\mathsf{send}}$、$t_{q,j}^{\mathsf{del}}$ 和 $t_{q,j}^{\mathsf{adopt}}$ 分别表示消息生成、发送、交付和采用时刻。

为排除数值噪声，令 $\varepsilon^{\mathsf u}>0$ 和 $\varepsilon^{\mathsf x}>0$ 分别为实际输入差异和状态差异的判定容差。从共同触发时刻 $t_0$ 开始，$t_j^{\mathsf u}$ 表示“存在源风险且开放消息”和“存在相同源风险但屏蔽消息”两种条件下，目标车辆的实际施加输入 $u_{j,t}$ 首次出现明显差异的时刻；$t_j^{\mathsf x}$ 表示在两种条件均屏蔽消息时，源风险存在与否首次使目标状态出现明显差异的时刻。物理影响到达时刻 $\tau_j^{\mathsf P}$ 沿用前文的时序物理路径审计定义，不用单个状态差异替代。$t_j^{\mathsf u}$ 用于信息先行审计，$t_j^{\mathsf x}$ 只用于校核物理传播估计。两个可观测分化时刻定义为
\begin{equation}
\begin{aligned}
t_j^{\mathsf u}&=\inf\left\{t\geq t_0:\left\|u_{j,t}(1,1)-u_{j,t}(1,0)\right\|>\varepsilon^{\mathsf u}\right\},\\
t_j^{\mathsf x}&=\inf\left\{t\geq t_0:\left\|x_{j,t}(1,0)-x_{j,t}(0,0)\right\|>\varepsilon^{\mathsf x}\right\},
\end{aligned}
\end{equation}
其中 $\inf$ 表示从 $t_0$ 开始寻找最早满足括号内条件的时刻；在离散仿真中，它对应第一个满足条件的时间步。差异还须连续保持预先冻结的时长，以排除单个时间步的偶然波动。若观察结束时仍未检测到分化，则约定 $\inf\varnothing=+\infty$，并在时间分析中记为右删失；在固定风险窗口的机制判定中，这种运行记为“尚未通过”，但仍保留在主效应分析中。对风险窗口 $\mathcal W=[t^{\mathsf{win}},t^{\mathsf{win}}+H)$，信息先行条件为
\begin{equation}
t_{q,j}^{\mathsf{adopt}}\leq t_j^{\mathsf u}<\tau_j^{\mathsf P},\qquad
t_{q,j}^{\mathsf{adopt}}\leq t^{\mathsf{win}}<t^{\mathsf{win}}+H\leq\tau_j^{\mathsf P}.
\end{equation}
为判断源风险能否通过普通交通交互到达目标车辆，令 $\operatorname{Reach}_{G^{\mathsf P}}(i,j;I)\in\{0,1\}$ 表示时间区间 $I$ 内是否存在从源车辆 $i$ 到目标车辆 $j$ 的时间单调物理路径。令 $\Sigma_{\omega}^{\mathsf{pre}}(s,c)$ 表示配对单元 $\omega$ 在处理条件 $(s,c)$ 下、处理开始前的冻结系统状态，符号 $\doteq_{\varepsilon^{\mathsf{pre}}}$ 表示离散字段与身份映射完全一致、连续状态之差不超过预先冻结的数值容差 $\varepsilon^{\mathsf{pre}}$。令 $e_\omega$ 为该单元指定的源风险事件，$\mathcal M_{\omega}^{\mathsf{src}}$ 为同时携带 $e_\omega$、源车辆 $i$、目标车辆 $j$、处理版本及有效载荷校验和的消息标识集合。令 $\mathcal K^{\mathsf{NC}}$ 表示预先指定的负对照集合，$\Delta_{k}^{\mathsf{NC}}$ 表示第 $k$ 类负对照产生的同类效应。Stage I 分别检验效应是否具有实际意义，以及该效应是否表现为安全收益；前者允许风险下降或增加，后者只允许风险下降。主问题及约束写为
\begin{equation}\label{eq:P1}
\begin{aligned}
\mathcal P_1:\quad
&\text{估计 }\Delta^{\mathsf{NL}}=\mathbb E_\omega[\delta_\omega^{\mathsf{NL}}];\\
&\mathsf{NLE}^{\mathsf{exist}}:\ |\Delta^{\mathsf{NL}}|\geq\delta^{\mathsf R};\\
&\mathsf{NLE}^{\mathsf{safe}}:\ \Delta^{\mathsf{NL}}\leq-\delta^{\mathsf R}.
\end{aligned}
\end{equation}
\addtocounter{equation}{-1}
下列约束中的 $q$ 表示满足约束~\eqref{eq:P1-d} 的一条已采用消息。
\begin{subequations}
\begin{align}
\mathrm{s.t.}\;&\mathcal R^{\mathsf S}\cap\mathcal R^{\mathsf T}=\varnothing;
\label{eq:P1-a}\\
&\operatorname{Reach}_{G^{\mathsf P}}\!(i,j;[t_0,t^{\mathsf{win}}+H])=0;
\label{eq:P1-b}\\
&\Sigma_{\omega}^{\mathsf{pre}}(s,c)\doteq_{\varepsilon^{\mathsf{pre}}}
\Sigma_{\omega}^{\mathsf{pre}}(s',c')
\quad\forall(s,c),(s',c')\in\{0,1\}^{2};
\label{eq:P1-c}\\
&\begin{aligned}
\exists q\in\mathcal M_{\omega}^{\mathsf{src}}:\quad
Z_{q,j}(t_{q,j}^{\mathsf{adopt}})&=1;
\end{aligned}
\label{eq:P1-d}\\
&Z_{q,j}(t)=1\Rightarrow D_{q,j}(t)=1\quad\forall t;
\label{eq:P1-e}\\
&Z_{q,j}(t)=1\Rightarrow Q_{q,j}(t)=1\quad\forall t;
\label{eq:P1-f}\\
&t_0\leq t_q^{\mathsf{gen}}\leq t_q^{\mathsf{send}};
\label{eq:P1-g}\\
&t_q^{\mathsf{send}}\leq t_{q,j}^{\mathsf{del}};
\label{eq:P1-h}\\
&t_{q,j}^{\mathsf{del}}\leq t_{q,j}^{\mathsf{adopt}};
\label{eq:P1-i}\\
&t_{q,j}^{\mathsf{adopt}}\leq t_j^{\mathsf u};
\label{eq:P1-j}\\
&t_j^{\mathsf u}<\tau_j^{\mathsf P};
\label{eq:P1-k}\\
&t_{q,j}^{\mathsf{adopt}}\leq t^{\mathsf{win}};
\label{eq:P1-l}\\
&t^{\mathsf{win}}+H\leq\tau_j^{\mathsf P};
\label{eq:P1-m}\\
&\left|\Delta_{k}^{\mathsf{NC}}\right|\leq\delta^{\mathsf{eq}}
\quad\forall k\in\mathcal K^{\mathsf{NC}}.
\label{eq:P1-n}
\end{align}
\end{subequations}
约束~\eqref{eq:P1-a}--\eqref{eq:P1-c} 是设计与处理前可比性条件。约束~\eqref{eq:P1-d} 将被归因的消息绑定到指定源事件、源车辆、目标车辆和处理版本；约束~\eqref{eq:P1-e}--\eqref{eq:P1-i} 审计该消息的交付、校验和生命周期顺序。约束~\eqref{eq:P1-j}--\eqref{eq:P1-m} 检查消息采用、实际动作和完整风险窗口是否先于物理到达边界。约束~\eqref{eq:P1-n} 在确认性实验层面分别检验错误来源、错误目标、错误类型、过期、仅交付未采用和虚假消息等负对照；每类负对照必须独立落入等效区间，不能先合并后判断。

令 $E_{\omega}^{\mathsf{tech}}=1$ 表示四格均按冻结配置完成，运行标识、校验和和主要结局可重建，且没有跨运行状态泄漏。运行完整性标签 $A_{\omega}^{\mathsf{run}}$ 与机制标签 $A_{\omega}^{\mathsf{mech}}$ 分别定义为
\begin{equation}
\begin{aligned}
A_{\omega}^{\mathsf{run}}=\mathsf{valid}
&\Longleftrightarrow E_{\omega}^{\mathsf{tech}}=1
\text{ 且约束 (9a)--(9c) 成立},\\
A_{\omega}^{\mathsf{mech}}=\mathsf{pass}
&\Longleftrightarrow\text{约束 (9d)--(9m) 成立}.
\end{aligned}
\end{equation}
其中 $A_{\omega}^{\mathsf{run}}\in\{\mathsf{valid},\mathsf{fail},\mathsf{cens}\}$ 记录工程完整性，$A_{\omega}^{\mathsf{mech}}\in\{\mathsf{pass},\mathsf{fail},\mathsf{cens}\}$ 记录作用链是否通过。全部预先登记的正式单元构成意向运行集合 $\Omega^{\mathsf{ITR}}$；$A_{\omega}^{\mathsf{mech}}$ 及任何处理后风险结果都不得用于移出该集合。工程失败和结局缺失保留在分母中，并按预先指定的缺失数据与界限分析处理。约束~\eqref{eq:P1-n} 在确认性实验层面检验负对照。只有效应界限、设计条件、机制审计和负对照均达到预设门槛时，才支持相应的 $\mathsf{NLE}^{\mathsf{exist}}$ 或 $\mathsf{NLE}^{\mathsf{safe}}$ 主张。

风险测量、物理路径和时间判定的解析真值、数值收敛和独立实现验证是 $\mathcal P_1$ 的前置有效性条件。Stage I 不优化消息采用策略或车辆控制策略，也不以消息到达、消息采用次数或单一风险指标的变化单独构成因果证据。

\subsection*{Stage II: safety-constrained control}
Stage II 仅在 $\mathcal P_1$ 的主要效应、消息链和负对照均达到预设门槛后启动。该启动条件不表示每条远端消息都必然有效，而是表示 Stage I 已经提供可重复证据，说明开放指定消息通道能够在预设条件下改变目标车辆的实际动作并改变目标区域总体风险。Stage II 随后回答控制器如何使用这类信息，并限制由此产生的安全、舒适性、效率和尾部代价。Stage II 的控制结果不能反过来证明 $\mathcal P_1$。

在物理连通场景中，令 $L_{\omega}^{\mathsf P}=\tau_j^{\mathsf P}-(t^{\mathsf{win}}+H)$ 表示完整风险窗口相对最早物理到达边界的领先裕度。若一条绑定指定源事件的消息在目标实际输入分化前被采用，且 $L_{\omega}^{\mathsf P}>0$，则该运行通过信息先行审计。没有通过的正式运行仍进入风险和控制结果的意向运行分析，只是不用于支持“信息先于物理影响”的归因。Result 4 只有在预先登记的全部归因运行均通过该审计，且物理到达估计的独立校准通过时，才报告信息先行结论。

Stage II 沿用消息交付、校验和采用状态 $D_{q,j}(t)$、$Q_{q,j}(t)$ 和 $Z_{q,j}(t)$。消息 $q$ 在时刻 $t$ 的年龄为 $t-t_q^{\mathsf{gen}}$，$T^{\mathsf{fresh}}>0$ 为最大允许年龄。令 $\widetilde M_{j,t}\subseteq M_{j,t}^{\mathsf{cand}}$ 为已经交付、通过联合校验且未过期的消息集合。采用策略 $\pi^{\mathsf Z}$ 为每条合格消息给出二元采用决定；被采用的消息组成 $M_{j,t}^{\mathsf{adopt}}$，再经冻结的消息聚合器形成控制信息 $z_{j,t}$。动作策略 $\pi^{\mathsf u}$ 根据本地状态和 $z_{j,t}$ 生成请求命令 $u_{j,t}^{\mathsf{cmd}}$，安全执行映射 $g_j^{\mathsf{safe}}$ 将其转化为实际输入 $u_{j,t}$。联合策略记为 $\pi=(\pi^{\mathsf Z},\pi^{\mathsf u})$。
具体地，合格消息集合定义为
\begin{equation*}
\begin{aligned}
\widetilde M_{j,t}=\{m_{q,i,j,t}\in M_{j,t}^{\mathsf{cand}}:
&D_{q,j}(t)=1,\ Q_{q,j}(t)=1,\\
&0\leq t-t_q^{\mathsf{gen}}\leq T^{\mathsf{fresh}}\}.
\end{aligned}
\end{equation*}

对运行 $\omega$，$J_\omega^{\mathsf{risk}}(\pi)$、$J_\omega^{\mathsf{act}}(\pi)$ 和 $J_\omega^{\mathsf{eff}}(\pi)$ 分别表示风险、动作与舒适性以及通行效率代价，且均按“代价越小越好”解释。令 $\mathbb P^{\mathsf{eval}}$ 为正式评估前冻结的运行分布。对 $\mathsf r\in\{\mathsf{risk},\mathsf{act},\mathsf{eff}\}$，总体代价和完整向量定义为
\begin{equation*}
\begin{aligned}
\mathbf J(\pi)&=\bigl(J^{\mathsf{risk}}(\pi),J^{\mathsf{act}}(\pi),
J^{\mathsf{eff}}(\pi),J^{\mathsf{tail}}(\pi)\bigr),\\
J^{\mathsf r}(\pi)&=\mathbb E_{\omega\sim\mathbb P^{\mathsf{eval}}}
\!\left[J_\omega^{\mathsf r}(\pi)\right].
\end{aligned}
\end{equation*}
这里先说明四类代价在 $\mathcal P_2$ 中的作用；其数学定义和缺失处理统一在后文的 Risk, action and temporal metrics 小节展开，尾部风险采用条件风险价值（CVaR）聚合。主问题最小化 $J^{\mathsf{risk}}(\pi)$，约束限制 $J^{\mathsf{act}}(\pi)$、$J^{\mathsf{eff}}(\pi)$ 和 $J^{\mathsf{tail}}(\pi)$。在一组预先冻结的代价上界上重复求解 $\mathcal P_2$，得到风险、安全、舒适性和效率之间的 Pareto 前沿，即无法继续改善某一目标而不使至少另一目标恶化的候选边界。最终策略选择规则在锁定留出集解锁前固定。令 $\mathcal U_j$ 为目标车辆允许的实际输入集合，$\mathcal X_j^{\mathsf{safe}}$ 为满足道路边界和碰撞规避要求的安全状态集合；$\overline J^{\mathsf{act}}$、$\overline J^{\mathsf{eff}}$ 和 $\overline J^{\mathsf{tail}}$ 为校准后冻结的代价上界。Stage II 的主问题及约束写为
\begin{equation}\label{eq:P2}
\begin{aligned}
\mathcal P_2:\quad
&\pi^\star\in\arg\min_{\pi}
J^{\mathsf{risk}}(\pi).
\end{aligned}
\end{equation}
\addtocounter{equation}{-1}
\allowdisplaybreaks[4]
\begin{subequations}
\begin{align}
\mathrm{s.t.}\;&x_{j,t+1}=f_j^{\mathsf{dyn}}(x_{j,t},u_{j,t},\xi_{j,t};
G^{\mathsf P}(t));
\label{eq:P2-a}\\
&u_{j,t}^{\mathsf{cmd}}=\pi^{\mathsf u}(x_{j,t},z_{j,t})
\quad\forall t;
\label{eq:P2-b}\\
&\begin{aligned}
Z_{q,j}(t)&=\pi_{q}^{\mathsf Z}(x_{j,t},m_{q,i,j,t})\in\{0,1\},\\
&\hspace{-2em}\forall m_{q,i,j,t}\in\widetilde M_{j,t};
\end{aligned}
\label{eq:P2-c}\\
&\begin{aligned}
M_{j,t}^{\mathsf{adopt}}=\{m_{q,i,j,t}\in\widetilde M_{j,t}:\\
Z_{q,j}(t)=1\};
\end{aligned}
\label{eq:P2-d}\\
&\begin{aligned}
z_{j,t}&=\operatorname{Agg}^{\mathsf{msg}}(M_{j,t}^{\mathsf{adopt}}),\\
\operatorname{Agg}^{\mathsf{msg}}(\varnothing)&=\mathbf 0;
\end{aligned}
\label{eq:P2-e}\\
&Z_{q,j}(t)=1\Rightarrow D_{q,j}(t)=1\quad\forall q,t;
\label{eq:P2-f}\\
&Z_{q,j}(t)=1\Rightarrow Q_{q,j}(t)=1\quad\forall q,t;
\label{eq:P2-g}\\
&Z_{q,j}(t)=1\Rightarrow
0\leq t-t_q^{\mathsf{gen}}\leq T^{\mathsf{fresh}}\quad\forall q,t;
\label{eq:P2-h}\\
&\begin{aligned}
u_{j,t}&=g_j^{\mathsf{safe}}(x_{j,t},u_{j,t}^{\mathsf{cmd}}),\\
u_{j,t}&\in\mathcal U_j\quad\forall t;
\end{aligned}
\label{eq:P2-i}\\
&x_{j,t}\in\mathcal X_j^{\mathsf{safe}}\quad\forall t;
\label{eq:P2-j}\\
&J^{\mathsf{act}}(\pi)\leq\overline J^{\mathsf{act}};
\label{eq:P2-k}
\\
&J^{\mathsf{eff}}(\pi)\leq\overline J^{\mathsf{eff}};
\label{eq:P2-l}\\
&J^{\mathsf{tail}}(\pi)\leq\overline J^{\mathsf{tail}}.
\label{eq:P2-m}
\end{align}
\end{subequations}
约束~\eqref{eq:P2-a} 使用实际输入更新车辆状态，约束~\eqref{eq:P2-b} 定义动作策略的请求命令。约束~\eqref{eq:P2-c}--\eqref{eq:P2-e} 依次定义逐消息采用、已采用消息集合和控制信息聚合，因此 $Z_{q,j}(t)$、$z_{j,t}$ 与控制命令具有可追溯的数学关系。约束~\eqref{eq:P2-f}--\eqref{eq:P2-h} 要求被采用的消息已经交付、通过校验且未过期。约束~\eqref{eq:P2-i} 将请求命令经过安全执行映射转化为实际输入，约束~\eqref{eq:P2-j} 要求全过程处于安全状态。约束~\eqref{eq:P2-k}--\eqref{eq:P2-m} 分别限制动作与舒适性、通行效率和尾部风险，防止策略通过急刹、停车或降低通行效率换取表面的平均风险改善。

控制器形式和参数在实验设计中冻结。仅本地控制器（local-only）是不接收远端消息的基线；固定保守控制器（always-cautious）不依赖消息并始终采取同一保守响应；理想信息控制器（oracle）使用完整先验风险信息，只作为不可部署的性能上界。候选策略只有在相对仅本地控制器取得具有实际意义的风险改善、相对固定保守控制器在安全和尾部风险上不劣，并在舒适性或效率至少一项上形成可验证改善时，才支持 Stage II 的控制价值主张。吞吐、排放、误报警和不必要采用作为预先指定的次要运行指标同时报告；它们不能替代四类核心代价或硬安全判定。若 $\mathcal P_2$ 未通过，只收缩控制价值主张，不推翻 $\mathcal P_1$ 的识别结论。

\subsection*{Experimental design and data sources}
独立重复单位是一个预先登记的“场景族 $\times$ 主随机种子”完整四格配对，而不是帧、消息、车辆时间步或同一运行中的多辆车。四格由同一初始快照启动，并使用按子系统命名且一一对应的共同随机流；它们只改变源风险指示 $S$ 和指定消息通道指示 $C$。场景族作为预先指定的分析层，不能在观察结果后合并或拆分。

$S=1$ 表示按冻结规则触发指定源风险事件，$S=0$ 表示在相同背景交通下不触发该事件。$C=1$ 表示指定消息协议栈开放：在 $S=1$ 时允许生成与该源事件绑定的风险消息，在 $S=0$ 时只运行等通信负载和等计算量的无风险占位消息，用于测量通道本身的副作用。$C=0$ 在协议入口屏蔽指定消息，但保留本地感知、基础控制器、安全约束和无关通信。车辆身份、道路规则和非处理控制逻辑在四格内保持一致。

正式运行前依次冻结校准集（calibration）、开发集（development）、确认集（confirmatory）和锁定留出集（locked holdout）的划分。校准集只确定测量规则和阈值，开发集可用于控制器调参，确认集只执行预注册的 Stage I 分析，锁定留出集在策略和分析规则固定后只解锁一次。场景清单、随机种子、车辆资格、源事件版本、消息协议、主要终点、效应边界、样本量和分析脚本同时冻结。

源事件只有在源车辆 $i$ 位于 $\mathcal R^{\mathsf S}$ 且满足冻结的风险触发规则时才合格；目标车辆 $j$ 必须在处理前基线和结果窗口内属于 $\mathcal R^{\mathsf T}$。物理隔离场景采用两条拓扑断开的走廊，并同时审计无跨走廊车辆转移、无共享控制状态和无时间单调物理路径；物理连通场景恢复预先指定的合流、交叉或换道关系，用于检验信息先行和控制迁移。每个场景在运行登记表中固定道路和车道范围、源区与目标区边界、拓扑距离、交通需求、车辆进入与退出规则、源扰动、背景驾驶模型、目标控制器、通信时延与丢包模型、仿真步长、观察期和风险窗口。上述内容只由校准和开发数据确定，不能根据确认性结果筛选配对或改写规则。

每次运行至少产生五类逻辑记录：车辆状态、源事件与标准化发射量、消息生命周期、控制器采用与请求/实际输入、风险和安全结果。每条记录带有运行标识、场景标识、随机种子、车辆或消息标识和仿真时间戳，并由运行清单（manifest）、来源记录（provenance）和失败账本关联。负对照包括错误来源、错误目标、错误风险类型、过期消息、仅交付未采用消息、等计算量的虚假消息和占位通道。真实轨迹与通信数据只用于其字段直接支持的测量、分布或可观测性校核\cite{ngsim2016dataset,coifman2017ngsim,barmpounakis2020pneuma,pneuma2024dataset,spmd2014dataset,zhan2019interaction}；缺少消息采用和实际动作字段时，不推断真实闭环因果效应。硬件在环或混合现实试验若用于外部闭环验证，必须采用相同的消息状态、时间同步、实际输入和风险字段，并作为独立证据层报告。

\subsection*{Risk, action and temporal metrics}
风险终点按冲突类型预先选择。追尾风险的确认性主终点使用归一化积分 TTC 负担；交叉、合流、换道和队列尾场景分别在校准阶段冻结适配的主要指标、方向、阈值、观测掩码、标准化和结果窗口。不同物理含义的原始 TTC、PET 或冲突量不直接相加，确认性结果也不能在解锁后更换主要指标。以下给出追尾 TTC 主终点的正式定义。令 $\varphi(j,t)$ 为时刻 $t$ 目标车辆 $j$ 的有效前车，$h_{j,t}>0$ 为净间距，$\Delta v_{j,t}=v_{j,t}-v_{\varphi(j,t),t}$ 为闭合速度；当不存在有效前车、间距或速度不可观测时，设观测掩码 $m_{j,t}^{\mathsf{TTC}}=0$，不以零替代。对有效观测，
\begin{equation*}
\operatorname{TTC}_{j,t}=
\begin{cases}
h_{j,t}/\Delta v_{j,t},&\Delta v_{j,t}>0,\\
+\infty,&\Delta v_{j,t}\leq 0.
\end{cases}
\end{equation*}

令 $\theta^{\mathsf{TTC}}>0$ 为校准后冻结的风险阈值，令 $[y]_+=\max(y,0)$。单个时间步的归一化风险负担为
\begin{equation*}
\rho_{j,t}^{\mathsf{TTC}}=
\begin{cases}
1,&\text{发生碰撞或 }\operatorname{TTC}_{j,t}=0,\\
1-\operatorname{TTC}_{j,t}/\theta^{\mathsf{TTC}},
&0<\operatorname{TTC}_{j,t}<\theta^{\mathsf{TTC}},\\
0,&\operatorname{TTC}_{j,t}\geq\theta^{\mathsf{TTC}}.
\end{cases}
\end{equation*}
在采用 TTC 主终点的运行中，$R_{j,\omega}(s,c)$ 就是对应处理格下的 $I_{j,\omega}^{\mathsf{TTC}}$；目标区域总体风险 $R_{\mathcal R^{\mathsf T},\omega}(s,c)$ 为目标车辆集合内单车风险负担之和。其他风险类型只在预先指定的指标替换后同步替换 $R$。若目标车辆 $j$ 的结果窗口内 $\sum_t m_{j,t}^{\mathsf{TTC}}\Delta t>0$，则
\begin{equation*}
\begin{aligned}
I_{j,\omega}^{\mathsf{TTC}}(\pi)
&:=\frac{\sum_{t\in\mathcal W_\omega}m_{j,t}^{\mathsf{TTC}}
\rho_{j,t}^{\mathsf{TTC}}\Delta t}
{\sum_{t\in\mathcal W_\omega}m_{j,t}^{\mathsf{TTC}}\Delta t},\\
J_{\omega}^{\mathsf{risk}}(\pi)
&:=\sum_{j\in\mathcal J^{\mathsf T}_{\omega}}I_{j,\omega}^{\mathsf{TTC}}(\pi).
\end{aligned}
\end{equation*}
无有效观测的运行记为不确定，不赋予零风险。碰撞、越界和不可恢复的安全事件同时作为硬安全结果记录，不因风险负担的平均值下降而被抵消。其他风险类型使用相同的“有效观测掩码—单位内归一化—预先冻结聚合”结构，并在跨类型比较前只比较标准化效应。

动作代价使用实际执行的加速度和加速度变化率（jerk），而不是控制器请求值。令 $m_{j,t}^{\mathsf a}\in\{0,1\}$ 表示二者在时刻 $t$ 是否同时可观测；当 $m_{j,t}^{\mathsf a}=0$ 时，该时间步不进入动作代价分母。令 $\bar a^{+}$ 和 $\bar a^{-}$ 为允许的加速和减速边界，$\bar a^{\mathsf{jerk}}$ 为 jerk 参考尺度，$e_{j,t}^{+}=[a_{j,t}-\bar a^{+}]_+$，$e_{j,t}^{-}=[-a_{j,t}-\bar a^{-}]_+$，并令 $\dot a_{j,t}=(a_{j,t}-a_{j,t-1})/\Delta t$。给定非负且和为一的权重 $w_{+}^{\mathsf a},w_{-}^{\mathsf a},w^{\mathsf{jerk}}$，时刻 $t$ 的动作代价为
\begin{equation*}
c_{j,t}^{\mathsf a}=
w_{+}^{\mathsf a}\left(\frac{e_{j,t}^{+}}{\bar a^{+}}\right)^2+
w_{-}^{\mathsf a}\left(\frac{e_{j,t}^{-}}{\bar a^{-}}\right)^2+
w^{\mathsf{jerk}}\left(\frac{\dot a_{j,t}}{\bar a^{\mathsf{jerk}}}\right)^2.
\end{equation*}
动作与舒适性总体代价为
\begin{equation*}
J_{\omega}^{\mathsf{act}}(\pi)
=\frac{\sum_{t\in\mathcal W_\omega}m_{j,t}^{\mathsf a}
c_{j,t}^{\mathsf a}\Delta t}
{\sum_{t\in\mathcal W_\omega}m_{j,t}^{\mathsf a}\Delta t}.
\end{equation*}
边界、参考尺度、权重和动作缺失规则在校准锁定前确定，确认性结果不得按策略表现重新调整。

效率代价记录旅行时间、停车时间和平均速度相对同一初始状态下参考运行的损失。令 $T_{\omega}^{\mathsf{obs}}(\pi)>0$ 为目标车辆完成规定任务所需的旅行时间；若观察期内未完成，则取冻结的最大观察时长 $T^{\mathsf{max}}$。令 $F_{\omega}^{\mathsf{inc}}(\pi)\in\{0,1\}$ 表示任务未完成，$T_{\omega}^{\mathsf{ref}}>0$ 为冻结参考控制器的旅行时间，$T_{\omega}^{\mathsf{stop}}(\pi)$ 为速度低于冻结停车阈值的时长，$\bar v_{\omega}(\pi)$ 和 $\bar v_{\omega}^{\mathsf{ref}}>0$ 为对应平均速度。给定非负且和为一的权重 $w^{\mathsf{time}},w^{\mathsf{stop}},w^{\mathsf{speed}},w^{\mathsf{inc}}$，
\begin{equation*}
\begin{aligned}
J_{\omega}^{\mathsf{eff}}(\pi)
={}&w^{\mathsf{time}}\left[
\frac{T_{\omega}^{\mathsf{obs}}(\pi)-T_{\omega}^{\mathsf{ref}}}
{T_{\omega}^{\mathsf{ref}}}\right]_+\\
&+w^{\mathsf{stop}}\frac{T_{\omega}^{\mathsf{stop}}(\pi)}
{T_{\omega}^{\mathsf{obs}}(\pi)}
+w^{\mathsf{speed}}\left[
1-\frac{\bar v_{\omega}(\pi)}{\bar v_{\omega}^{\mathsf{ref}}}\right]_+\\
&+w^{\mathsf{inc}}F_{\omega}^{\mathsf{inc}}(\pi).
\end{aligned}
\end{equation*}
参考运行、停车阈值、最大观察时长和权重只由校准集和开发集确定。候选策略未完成任务因此不能作为缺失而排除；若冻结参考运行本身不可重建，则整个配对记为工程失败并进入预设缺失数据界限分析。

尾部风险不以平均值代替。令 $\alpha\in(0,1)$ 为预先冻结的 CVaR 置信水平，因此被聚合的上尾概率为 $1-\alpha$；令 $\eta\in\mathbb R$ 为阈值变量。直接对运行级风险负担 $J_{\omega}^{\mathsf{risk}}(\pi)$ 做 CVaR 聚合，则总体尾部代价为
\begin{equation*}
J^{\mathsf{tail}}(\pi)=
\inf_{\eta\in\mathbb R}
\left\{\eta+\frac{1}{1-\alpha}
\mathbb E_{\omega\sim\mathbb P^{\mathsf{eval}}}\!\left[\left[J_{\omega}^{\mathsf{risk}}(\pi)-\eta\right]_+\right]\right\}.
\end{equation*}
四类代价均在同一冻结评估分布 $\mathbb P^{\mathsf{eval}}$ 上计算并按“越小越好”解释。$\alpha$、参考运行、权重、阈值和评估分布在正式评估前冻结。

时间先行判定使用消息采用时刻 $t_{q,j}^{\mathsf{adopt}}$、实际输入分化时刻 $t_j^{\mathsf u}$、物理到达边界 $\tau_j^{\mathsf P}$ 和风险窗口 $\mathcal W=[t^{\mathsf{win}},t^{\mathsf{win}}+H)$。$\inf$ 取离散时间中第一个满足容差并持续达到规定时长的时间步；观察结束仍未分化时保留为右删失。消息采用、实际输入分化和完整风险窗口均早于 $\tau_j^{\mathsf P}$ 时，该运行通过信息先行审计；未通过的正式运行保留在意向运行分母，并作为机制失败或不确定结果报告。

\subsection*{Statistical analysis and uncertainty}
四格配对单元 $\omega$ 是效应估计的最小分析单位。令 $\mathcal G$ 为预先登记的场景族集合，$\Omega_g\subseteq\Omega^{\mathsf{ITR}}$ 为场景族 $g$ 的全部正式配对单元，$n_g=|\Omega_g|$，$w_g^{\mathsf{scn}}\geq0$ 为冻结的场景权重且 $\sum_{g\in\mathcal G}w_g^{\mathsf{scn}}=1$。帧、消息、车辆时间步和同一运行内的多次事件只用于计算运行级结局。主要效应估计为
\begin{equation*}
\begin{aligned}
\widehat{\Delta}^{\mathsf{NL}}&=
\sum_{g\in\mathcal G}w_g^{\mathsf{scn}}
\left(\frac{1}{n_g}\sum_{\omega\in\Omega_g}
\widehat{\delta}_{\omega}^{\mathsf{NL}}\right),\\
\widehat{\delta}_{\omega}^{\mathsf{NL}}&=
\bigl[\widehat R_{\mathcal R^{\mathsf T},\omega}(1,1)-\widehat R_{\mathcal R^{\mathsf T},\omega}(1,0)\bigr]\\
&\quad-\bigl[\widehat R_{\mathcal R^{\mathsf T},\omega}(0,1)-\widehat R_{\mathcal R^{\mathsf T},\omega}(0,0)\bigr].
\end{aligned}
\end{equation*}
场景权重规定本文的目标总体是“各场景等权”还是“按预设交通暴露加权”，并在确认集解锁前固定。置信区间通过分层自助重采样（bootstrap）完整四格单元计算；若同一主随机种子跨场景复用，则共享该种子的单元作为一个簇共同重采样。若校准显示仍存在跨层依赖，则使用同一层级的混合效应模型作预先指定的敏感性分析，不改变目标效应。主要效应同时报告点估计、95\%置信区间、各场景分母、工程失败、机制失败和右删失数量。

确认性判定区分“存在实际重要效应”和“形成安全收益”。若 95\% 置信区间整体低于 $-\delta^{\mathsf R}$ 或整体高于 $+\delta^{\mathsf R}$，则支持 $\mathsf{NLE}^{\mathsf{exist}}$；只有预先指定的单侧 95\% 上界低于 $-\delta^{\mathsf R}$ 时，才支持 $\mathsf{NLE}^{\mathsf{safe}}$。正向越过 $+\delta^{\mathsf R}$ 表示实际重要的风险增加，必须作为有害效应报告。每类负对照效应还必须落入 $[-\delta^{\mathsf{eq}},\delta^{\mathsf{eq}}]$ 的等效区间。等效性使用双单侧检验或等价置信区间判据；不把“未达到显著性”当作等效。多个次要终点使用预先指定的 Holm 校正，其他比较明确标记为探索性。

缺失观测不作零值填补。因碰撞、越界、控制器中止或未完成任务造成的结局，按预先冻结的失败终点或最不利可实现界限计入，而不是删除。因日志损坏或基础设施故障无法重建主要结局时，该单元仍保留在 $\Omega^{\mathsf{ITR}}$ 的分母；主结果同时给出预先指定的缺失数据模型和保守界限分析，若结论依赖排除这些单元则确认性判定不通过。时间或路径在观察期内无法判断时记为右删失；只有预先指定的右删失模型经校准验证后才进行模型化处理。确认性样本量由校准集和开发集中的方差、配对相关性、$\delta^{\mathsf{eq}}$、$\delta^{\mathsf R}$ 和预设功效确定，并在确认集解锁前冻结。

Stage II 的策略选择只使用开发集和预先指定的稀有事件发现集；仅本地、固定保守和理想信息控制器在评估前冻结。候选策略必须满足 (11a)--(11m)，并相对仅本地控制器达到预设风险改善，相对固定保守控制器在安全和尾部风险上不劣，且在动作或效率至少一项上越过相应实际意义边界。Pareto 选择规则在锁定留出集解锁前固定，留出结果只用于一次性验证，不再调参。

\subsection*{Calibration, validation and reproducibility}
校准阶段只用于确定测量规则、容差、参考尺度、代价权重、CVaR 置信水平、效应边界、样本量和运行器参数，不用于估计确认性主效应。风险测量在人工构造的解析轨迹上验证 TTC、DRAC、PET 或相应冲突指标、碰撞判定和缺失掩码；真实轨迹校准集用于校核速度、间距、相对速度、减速度、事件持续时间和风险事件分布。极端减速度、长时低 TTC、轨迹边界和前车匹配错误进行分层人工审计。

时间展开物理路径算法使用独立实现或解析图例校验最早到达时刻，覆盖无路径、单路径、多路径、车辆离开、跨走廊隔离和右删失情形。消息状态机分别测试生成、发送、丢弃、延迟、交付、校验失败、采用和过期；控制器日志必须能由消息状态和车辆状态重建实际采用与动作输出。离散时间步 $\Delta t$ 通过数值收敛检查后固定，改变 $\Delta t$ 不得改变主要路径、消息链和结论方向。

每个运行保存代码提交标识、环境锁、配置和数据/随机种子 manifest、开始与结束时间、退出码、五类逻辑表、审计摘要、原始产物校验和及失败原因。修复 bug 后生成新版本并重跑全部受影响运行；不得用修复后的结果静默替换旧产物。分析脚本从只读原始表重建状态、消息生命周期、动作、风险和主估计量，并在结果解锁前通过一份独立实现或人工可审计样例复核。Results 报告验证结果、适用域和失败账本，Methods 只规定规则和复现条件。

\section*{Data availability}
本文生成的仿真状态、消息生命周期、控制动作、风险结果、Source Data、冻结配置和审计清单将在论文发表时随版本化研究归档公开。NGSIM、pNEUMA、Safety Pilot Model Deployment 和 INTERACTION 为第三方数据，分别受其原始访问条件和许可约束；本文不重新分发受限原始数据，但提供数据清单、校验和、处理代码及从合法取得的原始数据生成分析数据的步骤。

\section*{Code availability}
本研究的运行器、指标实现、统计分析、绘图脚本和复现说明维护于 \url{https://github.com/yys806/aad-traffic-risk-propagation}。论文对应版本将以固定发布标签、提交标识、环境锁和归档校验和标识；第三方仿真器与依赖软件按各自许可获取。

\begingroup
\sloppy
\raggedright
\Urlmuskip=0mu plus 1mu\relax
\begin{thebibliography}{99}

\bibitem{saberi2020contagion}
Saberi, M. \textit{et al.} A simple contagion process describes spreading of traffic jams in urban networks. \textit{Nat. Commun.} \textbf{11}, 1616 (2020). \url{https://doi.org/10.1038/s41467-020-15353-2}.

\bibitem{duan2023bottlenecks}
Duan, J. \textit{et al.} Spatiotemporal dynamics of traffic bottlenecks yields an early signal of heavy congestions. \textit{Nat. Commun.} \textbf{14}, 8002 (2023). \url{https://doi.org/10.1038/s41467-023-43591-7}.

\bibitem{wang2025delays}
Wang, J. \textit{et al.} Knowledge-guided self-learning control strategy for mixed vehicle platoons with delays. \textit{Nat. Commun.} \textbf{16}, 7705 (2025). \url{https://doi.org/10.1038/s41467-025-62597-x}.

\bibitem{wang2024conflict}
Wang, T. \textit{et al.} A conflict risk graph approach to modeling spatio-temporal dynamics of intersection safety. \textit{Transp. Res. C Emerg. Technol.} \textbf{169}, 104874 (2024). \url{https://doi.org/10.1016/j.trc.2024.104874}.

\bibitem{kolekar2020risk}
Kolekar, S., de Winter, J. \& Abbink, D. Human-like driving behaviour emerges from a risk-based driver model. \textit{Nat. Commun.} \textbf{11}, 4850 (2020). \url{https://doi.org/10.1038/s41467-020-18353-4}.

\bibitem{sun2020lookahead}
Sun, Y. \& Tan, C. On a class of new nonlocal traffic flow models with look-ahead rules. \textit{Physica D} \textbf{413}, 132663 (2020). \url{https://doi.org/10.1016/j.physd.2020.132663}.

\bibitem{huang2022nonlocal}
Huang, K. \& Du, Q. Stability of a nonlocal traffic flow model for connected vehicles. \textit{SIAM J. Appl. Math.} \textbf{82}, 221--243 (2022). \url{https://doi.org/10.1137/20M1355732}.

\bibitem{bao2022motifs}
Bao, X. \textit{et al.} Impact of basic network motifs on the collective response to perturbations. \textit{Nat. Commun.} \textbf{13}, 5301 (2022). \url{https://doi.org/10.1038/s41467-022-32913-w}.

\bibitem{kaiser2021isolators}
Kaiser, F., Latora, V. \& Witthaut, D. Network isolators inhibit failure spreading in complex networks. \textit{Nat. Commun.} \textbf{12}, 3143 (2021). \url{https://doi.org/10.1038/s41467-021-23292-9}.

\bibitem{hamedmoghadam2021percolation}
Hamedmoghadam, H., Jalili, M., Vu, H. L. \& Stone, L. Percolation of heterogeneous flows uncovers the bottlenecks of infrastructure networks. \textit{Nat. Commun.} \textbf{12}, 1254 (2021). \url{https://doi.org/10.1038/s41467-021-21483-y}.

\bibitem{luo2026crm}
Luo, Y., Li, X., Bhaskar, A., Ngoduy, D., Elhenawy, M. \& Glaser, S. Enhancing vehicle platoon safety assessment: A novel cascading risk measure. \textit{Accid. Anal. Prev.} (2026). \url{https://doi.org/10.1016/j.aap.2025.108360}.

\bibitem{razzaghpour2023controlaware}
Razzaghpour, M., Valiente, R., Zaman, M. \& Fallah, Y. P. Predictive model-based and control-aware communication strategies for cooperative adaptive cruise control. \textit{IEEE Open J. Intell. Transp. Syst.} \textbf{4}, 232--243 (2023). \url{https://doi.org/10.1109/OJITS.2023.3259283}.

\bibitem{wang2024topology}
Wang, X., Xu, C., Zhao, X., Li, H. \& Jiang, X. Stability and safety analysis of connected and automated vehicle platoon considering dynamic communication topology. \textit{IEEE Trans. Intell. Transp. Syst.} \textbf{25}, 13442--13452 (2024). \url{https://doi.org/10.1109/TITS.2024.3398111}.

\bibitem{ngsim2016dataset}
U.S. Department of Transportation Federal Highway Administration. Next Generation Simulation (NGSIM) vehicle trajectories and supporting data. Dataset (2016). \url{https://doi.org/10.21949/1504477}.

\bibitem{coifman2017ngsim}
Coifman, B. \& Li, L. A critical evaluation of the Next Generation Simulation (NGSIM) vehicle trajectory dataset. \textit{Transp. Res. B Methodol.} \textbf{105}, 362--377 (2017). \url{https://doi.org/10.1016/j.trb.2017.09.018}.

\bibitem{barmpounakis2020pneuma}
Barmpounakis, E. \& Geroliminis, N. On the new era of urban traffic monitoring with massive drone data: the pNEUMA large-scale field experiment. \textit{Transp. Res. C Emerg. Technol.} \textbf{111}, 50--71 (2020). \url{https://doi.org/10.1016/j.trc.2019.11.023}.

\bibitem{pneuma2024dataset}
Barmpounakis, E. \& Geroliminis, N. pNEUMA dataset. Zenodo (2024). \url{https://doi.org/10.5281/zenodo.10491409}.

\bibitem{spmd2014dataset}
U.S. Department of Transportation. Safety Pilot Model Deployment data. Dataset (2014). \url{https://doi.org/10.21949/1504482}.

\bibitem{zhan2019interaction}
Zhan, W. \textit{et al.} INTERACTION dataset: an international, adversarial and cooperative motion dataset in interactive driving scenarios with semantic maps. arXiv:1910.03088 (2019). \url{https://arxiv.org/abs/1910.03088}.

\bibitem{schumann2026collision}
Schumann, J. F. \textit{et al.} Active inference as a model of collision avoidance behavior in human drivers. \textit{Nat. Commun.} \textbf{17}, 5009 (2026). \url{https://doi.org/10.1038/s41467-026-73345-0}.

\bibitem{feng2026dense}
Feng, S. \textit{et al.} Breaking through safety performance stagnation in autonomous vehicles with dense learning. \textit{Nat. Commun.} \textbf{17}, 3163 (2026). \url{https://doi.org/10.1038/s41467-026-69761-x}.

\bibitem{xu2026ga2ad}
Xu, H. \textit{et al.} GA2AD: a generalized adaptive adversarial training framework considering surrounding hybrid risk field for autonomous driving. \textit{Nat. Commun.} (2026). \url{https://doi.org/10.1038/s41467-026-75643-z}.

\end{thebibliography}
\endgroup

\section*{Acknowledgements}
% 致谢与资助信息待写。

\section*{Author contributions}
% 作者贡献声明待写。

\section*{Competing interests}
% 利益冲突声明待写。

\section*{Figure legends}
% 图题与图注待写。

\section*{Tables}
% 表格待写。

\end{document}
